\documentclass[11pt,a4paper]{article}

% Packages
\usepackage[utf8]{inputenc}
\usepackage[T1]{fontenc}
\usepackage{amsmath,amssymb,amsthm}
\usepackage{hyperref}
\usepackage{booktabs}
\usepackage[margin=2.5cm]{geometry}

% Theorem environments
\newtheorem{theorem}{Theorem}[section]
\newtheorem{corollary}[theorem]{Corollary}
\newtheorem{definition}[theorem]{Definition}

\title{Structural Persistence Theory:\\Representation Theorem, Second Law, and Formal Verification}
\author{Akihito Sunagawa}
\date{}

\begin{document}
\maketitle

\begin{abstract}
We present Structural Persistence Theory (SPT), a framework for measuring how viable structure is irreversibly lost. The theory rests on two conditions: (1)~the viable set has positive measure, and (2)~repair is never free. A representation theorem, derived from the Cauchy functional equation under normalization, additivity, continuity, and nonnegativity, forces the stage-loss function to be $f(r) = -k \log r$. An impossibility theorem excludes all non-logarithmic alternatives. The persistence kernel $m(V_n) = m(V_0)\exp(-L_n)$ follows as a telescoping identity. With repair, the net loss $b_t = d_t - r_t$ replaces stage loss, giving $m(V_n) = m(V_0)\exp(-B_n)$. The structural second law---that cumulative total production $\Sigma_n$ is monotone nondecreasing---is proved as a necessary and sufficient characterization.

The formalization comprises 380 Lean~4 modules with zero sorry/admit/axiom. Conditional bridges to classical results---Shannon's uniqueness theorem, Jaynes' maximum entropy principle, Landauer's principle, the Crooks fluctuation theorem, and others---are constructed as accounting readouts under domain-specific witnesses. The theory's scope is characterized from both sides: every system satisfying the two conditions admits exactly one loss-measurement form, and every system violating either condition has a formally identified failure mode.
\end{abstract}

%----------------------------------------------------------------------
\section{Introduction}

\subsection{The question}

Structure can be lost even when resources remain. An organization may retain budget and personnel yet lose the ability to make coherent decisions. Software may retain computational resources yet become unmaintainable as dependency conflicts accumulate. A long-context language model may retain parameters yet lose logical consistency as unresolved contradictions build up.

In each case, what is lost is not the substrate but the set of states that can still sustain the structure. SPT formalizes this observation: structural loss is the shrinkage of the viable set, measured by a log-ratio scale that is uniquely forced by natural axioms.

\subsection{Contributions}

\begin{enumerate}
\item \textbf{Representation theorem}: Under normalization, additivity, continuity, and nonnegativity, the stage-loss function must be $f(r) = -k\log r$. No other form is consistent (\S\ref{sec:representation}).

\item \textbf{Structural second law}: Cumulative total production $\Sigma_n$ is monotone nondecreasing, proved as necessary and sufficient under two conditions: positive measure and non-free repair (\S\ref{sec:secondlaw}).

\item \textbf{Formal verification}: 380 Lean~4 modules, zero sorry/admit/axiom, with conditional bridges to 60+ fields (\S\ref{sec:lean}).
\end{enumerate}

\subsection{What this paper does not claim}

\begin{itemize}
\item SPT does not claim that all systems decay exponentially. The exponential form is a representation theorem for the viable-set measure, not an empirical law about any specific system.
\item SPT does not replace domain-specific dynamics. It provides the accounting coordinates; the dynamics are supplied by each domain.
\item The conditional bridges are not unconditional proofs of classical theorems. Each bridge requires domain-specific witnesses (choice of viable set~$V$, measure~$m$, and positivity verification).
\end{itemize}

%----------------------------------------------------------------------
\section{Setup}\label{sec:setup}

\subsection{Structural maintenance problem}

Fix a system~$X$. A \emph{structural maintenance problem} $\Pi = (X, G)$ consists of:
\begin{itemize}
\item A state space~$X$.
\item A maintenance condition~$G$ that defines which states can sustain the structure.
\item The \emph{viable set} $V_G = \{x \in X : x \text{ satisfies } G\}$, following Aubin's viability theory~\cite{aubin1991}.
\item A \emph{measure}~$m$ on~$X$, fixed before observation, that quantifies the ``room remaining.''
\end{itemize}

\subsection{Viable-set dynamics}

Constraints accumulate over time:
\[
V_0 \supseteq V_1 \supseteq \cdots \supseteq V_n.
\]
Each step consists of \emph{contraction} (constraints shrink~$V$) followed by \emph{repair} (recovery expands~$V$):
\[
V_t^{-} = K_t(V_t), \qquad V_{t+1} = R_t(V_t^{-}).
\]

\subsection{Stage loss and repair}

The \emph{stage loss} at step~$t$:
\[
d_t = -\log\frac{m(V_t^{-})}{m(V_t)}.
\]
The \emph{repair amount} at step~$t$:
\[
r_t = \log\frac{m(V_{t+1})}{m(V_t^{-})}.
\]
The \emph{net loss}:
\[
b_t = d_t - r_t.
\]

\subsection{Cumulative quantities}

Cumulative stage loss: $L_n = \sum_{t<n} d_t$. Cumulative net loss: $B_n = \sum_{t<n} b_t$. Effective resource: $M_n$ (resource-side scalar, defined separately).

\subsection{Total production}

Under a \emph{repair budget} (repair gain~$\leq$ repair cost at each step), total production is:
\[
\Sigma_n = B_n + C_n = L_n + \text{slack}_n,
\]
where $C_n$ is cumulative repair cost and $\text{slack}_n = C_n - (\text{cumulative gain}) \geq 0$.

%----------------------------------------------------------------------
\section{The Representation Theorem}\label{sec:representation}

\subsection{Axioms for stage loss}

A \emph{stage-loss functional} $f \colon (0,1] \to \mathbb{R}$ satisfying:
\begin{itemize}
\item[\textbf{B2}] (Normalization): $f(1) = 0$.
\item[\textbf{B3}] (Additivity): $f(r_1 r_2) = f(r_1) + f(r_2)$.
\item[\textbf{B4}] (Continuity): $f$ is continuous.
\item[\textbf{NN}] (Nonnegativity): $f(r) \geq 0$ for $r \in (0,1]$.
\end{itemize}

\subsection{Uniqueness}

\begin{theorem}[Representation]\label{thm:representation}
Any stage-loss functional satisfying B2+B3+B4+NN must have
\[
f(r) = -k\log r, \quad k \geq 0.
\]
\end{theorem}

\begin{proof}[Proof sketch]
B3 is the Cauchy functional equation $f(xy) = f(x) + f(y)$. Under B4 (continuity), the unique solution is $f(r) = c \cdot \log r$. B2 is automatic (derivable from B3 alone). Nonnegativity on $(0,1]$ forces $c \leq 0$, giving $f(r) = -k\log r$ with $k \geq 0$. Lean: \texttt{RepresentationTheorem.loss\_must\_be\_log}.\qed
\end{proof}

\begin{corollary}[Persistence kernel]
For any positive mass sequence:
\[
m(V_n) = m(V_0) \cdot \exp(-L_n),
\]
where $L_n = \sum l_i$. At $k=1$ (structural nats), $S = M\exp(-L)$. Lean: \texttt{TelescopingExp.measure\_eq\_initial\_mul\_exp\_neg\_cumulative\_loss}.
\end{corollary}

\subsection{Impossibility}

\begin{theorem}[Impossibility]\label{thm:impossibility}
No non-logarithmic function satisfies B2+B3+B4+NN. Specifically:
\begin{itemize}
\item Linear $a(1-r)$ violates B3.
\item Quadratic $a(1-r)^2$ violates B3.
\end{itemize}
Lean: \texttt{ImpossibilityTheorem.impossibility\_of\_non\_log}.
\end{theorem}

\subsection{Relation to Shannon}

The representation theorem shares the same mathematical skeleton as Shannon's uniqueness theorem for entropy~\cite{shannon1948}. Both derive from the Cauchy functional equation; both force a logarithmic form. The difference is in the domain: Shannon measures message uncertainty; SPT measures viable-set shrinkage. Lean: \texttt{ShannonUniquenessCorollary}. SPT does not claim to be Shannon's theorem; it shares a common mathematical ancestor and records the correspondence (\texttt{NonIdentityTheorem}).

%----------------------------------------------------------------------
\section{The Structural Second Law}\label{sec:secondlaw}

\subsection{Statement}

\begin{theorem}[Structural Second Law]\label{thm:secondlaw}
Under:
\begin{enumerate}
\item Positive trajectory: $m(V_t) > 0$ for all~$t$.
\item Repair budget: gain~$\leq$ cost at each step.
\end{enumerate}
Cumulative total production $\Sigma_n$ is monotone nondecreasing. Lean: \texttt{StructuralSecondLaw.deterministic\_second\_law}.
\end{theorem}

\begin{theorem}[Converse]\label{thm:converse}
$\Sigma_n$ monotone $\iff$ $\text{stepTotalProduction} \geq 0$ at every step. Lean: \texttt{ConverseSecondLaw}.
\end{theorem}

\subsection{Three layers}

\begin{center}
\begin{tabular}{lll}
\toprule
Layer & Statement & Lean module \\
\midrule
Deterministic & $\Sigma_n$ pointwise monotone & \texttt{deterministic\_second\_law} \\
Stochastic & $\mathbb{E}[\Sigma_n]$ monotone & \texttt{stochastic\_second\_law} \\
Coarse-graining & Monotonicity transfers & \texttt{coarse\_second\_law} \\
\bottomrule
\end{tabular}
\end{center}

\subsection{Necessity of conditions}

\begin{theorem}[Free Repair Impossibility]\label{thm:freerepair}
If gain $>$ cost at any step, $\Sigma$ can decrease. The repair budget is necessary. Lean: \texttt{FreeRepairImpossibility}.
\end{theorem}

\begin{theorem}[Minimal Axioms]\label{thm:minimal}
The two conditions are minimal and jointly sufficient. Lean: \texttt{MinimalAxiomTheorem}.
\end{theorem}

%----------------------------------------------------------------------
\section{Complete Scope Closure}

\subsection{Applicable systems}

Every system with $m > 0$ and gain~$\leq$ cost admits the unique form $S = M\exp(-L)$. The two conditions are necessary and sufficient. The functional form is unique and the axiom system is minimal.

\subsection{Inapplicable systems}

\begin{center}
\begin{tabular}{lll}
\toprule
Condition violated & Failure mode & Lean proof \\
\midrule
$m = 0$ & Log-ratio undefined & \texttt{ScopeBoundaryTheorem} \\
$m < 0$ & Physically meaningless & \texttt{CompleteScopeClosure} \\
gain $>$ cost & Second law violated & \texttt{FreeRepairImpossibility} \\
Constant mass & $L = 0$ (theory trivial) & \texttt{ScopeBoundaryTheorem} \\
\bottomrule
\end{tabular}
\end{center}

No fifth category exists. Lean: \texttt{CompleteScopeClosure.classification\_exhaustive}.

%----------------------------------------------------------------------
\section{Lean 4 Formalization}\label{sec:lean}

\subsection{Scale}

380 modules, 3{,}500+ build jobs. sorry~$= 0$, admit~$= 0$, axiom~$= 0$. Lean~4 v4.26.0 + Mathlib v4.26.0. Repository: \url{https://github.com/karesansui-u/persistence-lean}.

\subsection{Architecture}

\begin{center}
\begin{tabular}{lll}
\toprule
Layer & Modules & Content \\
\midrule
0--4 & Core kernel & Telescoping, log-uniqueness, resource budget \\
5--6 & SAT/CSP & Finite-horizon Chernoff/KL bounds \\
7--8 & Markov & Foster-Lyapunov, repair chains \\
9 & Route A & Serial reliability, BSC, fatigue \\
10--11 & Unification & CrossClass, StructuralSecondLaw \\
Tier S & Meta-theorems & Representation, Impossibility \\
Tier S+ & Necessity & FreeRepair, MinimalAxiom, Converse \\
Tier S++ & Completeness & Stability, Duality, Invariance \\
Tier S+++ & Validation & Falsifiability, NonIdentity \\
\bottomrule
\end{tabular}
\end{center}

\subsection{Conditional bridges}

Each bridge formalizes the algebraic readout of the SPT kernel in a specific domain. These are not unconditional proofs of classical theorems but conditional accounting correspondences under domain-specific witnesses. Bridges vary in depth: Tier~A bridges invoke the core (representation theorem, telescoping kernel, or second law) to derive non-trivial domain conclusions; Tier~B bridges provide correct structural correspondences; Tier~C bridges record vocabulary mappings. All are sorry-free, but they are not all at the same depth.

60+ fields are connected, including: thermodynamics, quantum mechanics, statistical mechanics, information theory, probability, biology, economics, engineering, and cosmology.

%----------------------------------------------------------------------
\section{Three Observability Layers}

\begin{center}
\begin{tabular}{lll}
\toprule
Layer & Observability & Role \\
\midrule
Specification-fixed & $V$, $m$, boundary fixed & Theorems, bounds \\
Conditional embedding & Drift mapped to SPT & Formal bridges \\
Structural estimation & Proxy indicators & Prediction, diagnosis \\
\bottomrule
\end{tabular}
\end{center}

These are not three different theories. They are three observability levels of the same kernel $S = M\exp(-L)$.

%----------------------------------------------------------------------
\section{Related Work}

SPT relates to but is structurally distinct from:
\begin{itemize}
\item \textbf{Thermodynamics}: SPT allows $B_n < 0$ (net repair exceeds loss); Clausius entropy is nondecreasing in isolated systems. SPT models open systems with repair.
\item \textbf{Information theory}: SPT measures viable-set shrinkage, not message uncertainty. Same functional form, different domain~\cite{shannon1948}.
\item \textbf{Survival analysis}: SPT tracks set-valued dynamics with repair; $S = M\exp(-L)$ can exceed~1.
\item \textbf{Viability theory}~\cite{aubin1991}: SPT adds an accounting layer---measuring how much the viable kernel has shrunk, not just whether viable trajectories exist.
\end{itemize}

%----------------------------------------------------------------------
\section{Limitations}

\begin{itemize}
\item \textbf{No dynamics}: SPT accounts for structural loss but does not supply dynamical equations.
\item \textbf{$G$ is external}: The choice of maintenance condition requires domain knowledge.
\item \textbf{Finite horizon}: Core results are finite-horizon. Asymptotic extensions require additional assumptions.
\item \textbf{Bridges are conditional}: Each domain bridge requires domain-specific witnesses.
\end{itemize}

%----------------------------------------------------------------------
\section{Conclusion}

Structural Persistence Theory provides a formally verified framework for measuring irreversible structural loss. Two conditions---positive measure and non-free repair---are necessary and sufficient. The representation theorem forces the unique form $f(r) = -k\log r$. The impossibility theorem excludes alternatives. The complete scope closure characterizes applicability from both sides.

The breadth of conditional bridges follows from mathematical structure: the Cauchy functional equation admits only one continuous solution, and many domains satisfy the two conditions under appropriate witnesses. This breadth is not a claim of universality by fiat, but a consequence of the axioms being weak enough to apply broadly while still forcing a unique functional form.

The Lean~4 formalization with 380 modules and zero sorry/admit provides machine-checked confidence in the mathematical core. The domain bridges provide conditional correspondences, each requiring explicit witnesses, and each open to independent verification.

%----------------------------------------------------------------------
\begin{thebibliography}{13}
\bibitem{aubin1991} Aubin, J.-P. (1991). \emph{Viability Theory}. Birkh\"auser.
\bibitem{aubin2011} Aubin, J.-P., Bayen, A.M., and Saint-Pierre, P. (2011). \emph{Viability Theory: New Directions}. Springer.
\bibitem{shannon1948} Shannon, C.E. (1948). A mathematical theory of communication. \emph{Bell System Technical Journal}, 27, 379--423.
\bibitem{khinchin1957} Khinchin, A.Ya. (1957). \emph{Mathematical Foundations of Information Theory}. Dover.
\bibitem{crooks1999} Crooks, G.E. (1999). Entropy production fluctuation theorem. \emph{Physical Review E}, 60, 2721.
\bibitem{jarzynski1997} Jarzynski, C. (1997). Nonequilibrium equality for free energy differences. \emph{Physical Review Letters}, 78, 2690.
\bibitem{friston2010} Friston, K. (2010). The free-energy principle: a unified brain theory? \emph{Nature Reviews Neuroscience}, 11, 127--138.
\bibitem{landauer1961} Landauer, R. (1961). Irreversibility and heat generation in the computing process. \emph{IBM Journal of Research and Development}, 5, 183--191.
\bibitem{fisher1930} Fisher, R.A. (1930). \emph{The Genetical Theory of Natural Selection}. Clarendon Press.
\bibitem{boltzmann1877} Boltzmann, L. (1877). \"Uber die Beziehung zwischen dem zweiten Hauptsatze der mechanischen W\"armetheorie und der Wahrscheinlichkeitsrechnung. \emph{Wiener Berichte}, 76, 373--435.
\bibitem{doob1953} Doob, J.L. (1953). \emph{Stochastic Processes}. Wiley.
\bibitem{hyers1941} Hyers, D.H. (1941). On the stability of the linear functional equation. \emph{PNAS}, 27, 222--224.
\bibitem{banach1922} Banach, S. (1922). Sur les op\'erations dans les ensembles abstraits. \emph{Fundamenta Mathematicae}, 3, 133--181.
\end{thebibliography}

\appendix
\section{Repository}

Lean formalization: \url{https://github.com/karesansui-u/persistence-lean}

Paper repository: \url{https://github.com/karesansui-u/delta-survival-papers}

Build: \texttt{lake build Persistence} (Lean~4 v4.26.0 + Mathlib v4.26.0)

380 modules. 3{,}500+ build jobs. sorry~$= 0$. admit~$= 0$. axiom~$= 0$.

\end{document}
